\documentclass[11pt,a4paper]{article}
\usepackage[utf8]{inputenc}
\usepackage[T1]{fontenc}
\usepackage{amsmath,amsfonts,amssymb}
\usepackage{graphicx}
\usepackage{hyperref}
\usepackage{listings}
\usepackage{xcolor}
\usepackage{booktabs}
\usepackage{float}
\usepackage{geometry}
\geometry{margin=1in}

% Color definitions
\definecolor{zenblue}{RGB}{41,121,255}
\definecolor{zengreen}{RGB}{52,199,89}
\definecolor{zenorange}{RGB}{255,149,0}
\definecolor{codegray}{RGB}{245,245,245}

% Hyperref setup
\hypersetup{
    colorlinks=true,
    linkcolor=zenblue,
    urlcolor=zenblue,
    citecolor=zenblue
}

% Code listing setup
\lstset{
    backgroundcolor=\color{codegray},
    basicstyle=\ttfamily\small,
    breaklines=true,
    captionpos=b,
    frame=single,
    numbers=left,
    numberstyle=\tiny\color{gray}
}

\title{
    \vspace{-2cm}
    \Large \textbf{Zen AI Model Family} \\
    \vspace{0.5cm}
    \Huge \textbf{Zen-Scribe} \\
    \vspace{0.3cm}
    \large A Packaged Multilingual Speech-Recognition Distribution \\
    built on Qwen3-ASR \\
    \vspace{0.5cm}
    \normalsize Technical Whitepaper v1.0
}

\author{
    Zach Kelling\thanks{zach@lux.network} \\
    \texttt{research@hanzo.ai} \\
    \\
    Zoo Labs Foundation \\
    \texttt{foundation@zoolabs.org}
}

\date{September 2025}

\begin{document}

\maketitle

\begin{abstract}
\textbf{Zen-Scribe} is a packaged, deployment-ready distribution of Alibaba's
open-source \textbf{Qwen3-ASR} speech-recognition model~\cite{qwen3asr}. It is
\emph{not} a model trained from scratch: the recognition capability, multilingual
coverage, and accuracy are entirely those of the upstream Qwen3-ASR weights, which
are released under the Apache~2.0 license. Zen-Scribe contributes packaging, not new
training---quantized build artifacts (GGUF, MLX), a thin transcription API, and
documentation---so that the upstream model is easy to self-host on commodity and
Apple-Silicon hardware. This whitepaper documents what Zen-Scribe wraps, the
provenance of the underlying model, and how to deploy it. All capability claims are
attributed to the upstream Qwen3-ASR technical report rather than re-measured here.
\end{abstract}

\tableofcontents
\newpage

\section{Introduction}

Self-hosting a strong multilingual speech-recognition model is often blocked not by
model quality but by packaging: converting weights to quantized runtime formats,
wiring a transcription interface, and documenting hardware requirements.
\textbf{Zen-Scribe} addresses this gap by repackaging an existing, openly-licensed
ASR model---Alibaba's \textbf{Qwen3-ASR}~\cite{qwen3asr}---for convenient local
deployment. We make no claim to have trained a recognition model; our contribution is
distribution and tooling around upstream weights.

\subsection{What Zen-Scribe Provides}
\begin{itemize}
    \item \textbf{Packaging}: GGUF and MLX build artifacts of the upstream Qwen3-ASR
          weights for CPU, GPU, and Apple-Silicon inference.
    \item \textbf{A thin API}: A \texttt{transcribe()} convenience wrapper over the
          upstream model.
    \item \textbf{Documentation}: Hardware sizing and deployment notes (below).
\end{itemize}

\subsection{What Zen-Scribe Does Not Provide}
\begin{itemize}
    \item No new pretraining, fine-tuning, or distillation of the acoustic model.
    \item No independently-measured accuracy numbers (see \S\ref{sec:upstream}).
    \item No proprietary architecture; the architecture is Qwen3-ASR's.
\end{itemize}

\section{Underlying Model}
\label{sec:upstream}

\subsection{Provenance}

Zen-Scribe wraps \textbf{Qwen3-ASR}, the open-source automatic-speech-recognition
series released by the Qwen team at Alibaba Cloud
(\texttt{Qwen3ASRForConditionalGeneration})~\cite{qwen3asr}. The upstream family is
released under the \textbf{Apache~2.0} license, which permits redistribution and
commercial use; Zen-Scribe inherits that license and redistributes the weights
unmodified.

\begin{table}[H]
\centering
\begin{tabular}{ll}
\toprule
\textbf{Component} & \textbf{Specification} \\
\midrule
Upstream model & Qwen3-ASR (Alibaba Cloud / Qwen team) \\
Architecture & As published upstream (encoder + LLM decoder) \\
HF class & \texttt{Qwen3ASRForConditionalGeneration} \\
License & Apache 2.0 (inherited from upstream) \\
Capabilities & Multilingual ASR, language ID, timestamps \\
\bottomrule
\end{tabular}
\caption{Underlying model provenance. All specifications are inherited from Qwen3-ASR.}
\end{table}

\subsection{Capabilities (as reported upstream)}

The following capabilities are those documented by the upstream Qwen3-ASR technical
report~\cite{qwen3asr}; we cite rather than re-measure them. Per Alibaba's release,
the Qwen3-ASR series supports stable multilingual speech, music, and song recognition
across a large set of languages and dialects, with automatic language detection,
streaming and offline operation from a unified model, and timestamp prediction. For
exact word-error-rate figures, benchmark methodology, and per-language breakdowns,
readers should consult the upstream report directly; we deliberately do not reproduce
benchmark tables we did not run.

\subsection{``Flash'' variant}

Alibaba additionally offers a hosted \textbf{Qwen3-ASR-Flash} real-time WebSocket
endpoint via the DashScope platform. That hosted service is distinct from the
self-hostable open weights that Zen-Scribe packages; Zen-Scribe is a local
distribution and does not proxy the hosted API.

\section{Packaging and Tooling}

Zen-Scribe's engineering work is confined to packaging:
\begin{enumerate}
    \item \textbf{Format conversion}: Producing GGUF (\texttt{Q4\_K\_M},
          \texttt{Q5\_K\_M}, \texttt{Q8\_0}) and MLX (4-bit, 8-bit) artifacts from
          the upstream Apache-2.0 weights.
    \item \textbf{Interface}: A thin \texttt{transcribe()} wrapper.
    \item \textbf{Verification}: Smoke tests confirming the packaged artifacts produce
          output equivalent to the upstream reference implementation on sample audio.
\end{enumerate}

We do not report custom accuracy or throughput benchmarks here, because the model's
behavior is that of upstream Qwen3-ASR; any independently-run numbers would need their
own documented methodology before they could be cited.

\section{Use Cases and Applications}

\subsection{Primary Applications}
\begin{itemize}
\item Real-time transcription
\item Meeting notes and summaries
\item Podcast transcription
\item Multilingual subtitles
\item Voice command processing
\end{itemize}

\subsection{Integration Examples}

\begin{lstlisting}[language=Python, caption=Basic Usage Example (loads the upstream Qwen3-ASR weights repackaged by Zen-Scribe)]
import librosa
from transformers import AutoProcessor, AutoModelForSpeechSeq2Seq

# Load the packaged upstream Qwen3-ASR model
model = AutoModelForSpeechSeq2Seq.from_pretrained("zenlm/zen-scribe")
processor = AutoProcessor.from_pretrained("zenlm/zen-scribe")

# Transcribe
audio, sr = librosa.load("speech.wav", sr=16000)
inputs = processor(audio, sampling_rate=sr, return_tensors="pt")
transcription = processor.batch_decode(
    model.generate(**inputs), skip_special_tokens=True
)
print(transcription[0])
\end{lstlisting}

\section{Safety, Privacy, and Responsible Use}

Because Zen-Scribe redistributes upstream weights unmodified, its model-level safety
properties are those of Qwen3-ASR; we add no new alignment training. At the
deployment layer we recommend:
\begin{itemize}
    \item \textbf{Privacy}: Transcription can run fully on-device; audio need not leave
          the host. Operators should obtain consent before transcribing third-party
          speech and avoid retaining raw audio beyond what is required.
    \item \textbf{Limitations}: ASR accuracy varies by language, accent, audio quality,
          and domain. Per-language behavior is inherited from upstream; consult the
          Qwen3-ASR report~\cite{qwen3asr} for known limitations.
    \item \textbf{Bias}: Recognition error rates can differ across accents and
          dialects; downstream applications should account for this.
\end{itemize}

\section{Deployment Options}

\subsection{Available Formats}
\begin{itemize}
    \item \textbf{SafeTensors}: Upstream-precision weights (as released by Alibaba).
    \item \textbf{GGUF}: Quantized formats (Q4\_K\_M, Q5\_K\_M, Q8\_0).
    \item \textbf{MLX}: Apple Silicon optimization (4-bit, 8-bit).
\end{itemize}

\subsection{Hardware Requirements}

Memory footprint depends on the upstream parameter count and the chosen quantization.
As a rule of thumb, lower-precision quantization reduces memory monotonically (Q4
$<$ Q8 $<$ FP16); the values below are approximate sizing guidance for a model in the
small-LLM range and should be validated against the specific upstream checkpoint and
runtime.

\begin{table}[H]
\centering
\begin{tabular}{lll}
\toprule
\textbf{Precision} & \textbf{Approx. Memory} & \textbf{Example Hardware} \\
\midrule
FP16 & higher & Discrete GPU (e.g.\ RTX 3060+) \\
INT8 & medium & Entry GPU / high-RAM CPU \\
INT4 & lowest & CPU / small-footprint devices \\
\bottomrule
\end{tabular}
\caption{Approximate hardware sizing by precision. Exact memory follows from the
upstream checkpoint size; measure before provisioning.}
\end{table}

\section{Future Work}

\begin{itemize}
    \item Tracking upstream Qwen3-ASR releases and re-packaging new checkpoints.
    \item Additional runtime targets (e.g.\ ONNX) for the same upstream weights.
    \item Optional, clearly-labelled fine-tunes for specific domains, kept separate
          from the base redistribution.
\end{itemize}

\section{Conclusion}

\textbf{Zen-Scribe} is a packaging and deployment layer over Alibaba's openly-licensed
\textbf{Qwen3-ASR} model. Its value is convenience---quantized artifacts, a thin API,
and documentation---not novel modeling. The recognition quality is entirely upstream's;
we attribute it accordingly and avoid presenting benchmarks we did not run.

\section*{Acknowledgments}

We thank the Qwen team at Alibaba Cloud for releasing Qwen3-ASR under Apache~2.0, and
the broader open-source community whose tooling makes local repackaging possible.

\begin{thebibliography}{99}
\bibitem{qwen3asr} Qwen Team, Alibaba Cloud. (2026). Qwen3-ASR Technical Report. arXiv:2601.21337. Models and code: \url{https://github.com/QwenLM/Qwen3-ASR} (Apache 2.0).
\bibitem{vaswani2017attention} Vaswani, A. et al. (2017). Attention Is All You Need. NeurIPS 2017.
\bibitem{brown2020language} Brown, T. et al. (2020). Language Models are Few-Shot Learners. NeurIPS 2020.
\bibitem{radford2019language} Radford, A. et al. (2019). Language Models are Unsupervised Multitask Learners. OpenAI Technical Report.
\bibitem{ouyang2022training} Ouyang, L. et al. (2022). Training Language Models to Follow Instructions with Human Feedback. NeurIPS 2022.
\bibitem{raffel2020exploring} Raffel, C. et al. (2020). Exploring the Limits of Transfer Learning with a Unified Text-to-Text Transformer. JMLR 2020.
\bibitem{radford2023robust} Radford, A. et al. (2023). Robust Speech Recognition via Large-Scale Weak Supervision. ICML 2023.
\bibitem{schulman2017proximal} Schulman, J. et al. (2017). Proximal Policy Optimization Algorithms. arXiv:1707.06347.
\bibitem{rafailov2023direct} Rafailov, R. et al. (2023). Direct Preference Optimization: Your Language Model is Secretly a Reward Model. NeurIPS 2023.
\bibitem{touvron2023llama} Touvron, H. et al. (2023). LLaMA: Open and Efficient Foundation Language Models. arXiv:2302.13971.
\bibitem{hu2021lora} Hu, E. et al. (2021). LoRA: Low-Rank Adaptation of Large Language Models. ICLR 2022.
\bibitem{dettmers2023qlora} Dettmers, T. et al. (2023). QLoRA: Efficient Finetuning of Quantized Language Models. NeurIPS 2023.
\end{thebibliography}

\appendix

\section{Model Card}

\begin{table}[H]
\centering
\begin{tabular}{ll}
\toprule
\textbf{Field} & \textbf{Value} \\
\midrule
Distribution Name & Zen-Scribe \\
Upstream Model & Qwen3-ASR (Alibaba Cloud / Qwen team) \\
Relationship & Repackaging of upstream weights (no retraining) \\
Version & 1.0.0 \\
License & Apache 2.0 (inherited from upstream) \\
Repository & \href{https://huggingface.co/zenlm/zen-scribe}{huggingface.co/zenlm/zen-scribe} \\
Upstream & \href{https://github.com/QwenLM/Qwen3-ASR}{github.com/QwenLM/Qwen3-ASR} \\
Contact & research@hanzo.ai \\
\bottomrule
\end{tabular}
\caption{Distribution card. Zen-Scribe is a packaged distribution of Qwen3-ASR, not an
independently trained model.}
\end{table}

\end{document}